\documentclass[11pt,letterpaper]{article}

% Margins and spacing (JFE-ish: narrow-ish margins, 1.2 line spacing)
\usepackage[letterpaper,margin=1in]{geometry}
\usepackage{setspace}
\onehalfspacing

% Core packages
\usepackage[utf8]{inputenc}
\usepackage[T1]{fontenc}
\usepackage{lmodern}
\usepackage{amsmath,amssymb}
\usepackage{booktabs}
\usepackage{longtable}
\usepackage{array}
\usepackage[round,authoryear]{natbib}
\usepackage{graphicx}
\usepackage[colorlinks=true,linkcolor=black,citecolor=black,urlcolor=blue]{hyperref}
\usepackage{url}
\usepackage{xcolor}

% Custom macro for TODO markers the user should resolve before SSRN post
\newcommand{\TODO}[1]{\textcolor{red}{\textbf{[TODO: #1]}}}

\title{The Della Vedova Prediction Market Indices (DV-PMI):\\
       Construction, Validation, and Use}

\author{Joshua Della Vedova%
  \thanks{Knauss School of Business, University of San Diego,
    5998 Alcala Park, San Diego, CA 92110.
    Email: \texttt{jdellavedova@sandiego.edu}.
    Dashboard and data: \url{https://jdellavedova.com}.
    Source code: \url{https://github.com/jdellavedova/polymarket-index}.}
}

\date{\today}

\begin{document}

\maketitle

\begin{abstract}
\noindent
Prediction markets with binary terminal payoffs allow direct measurement of trader behavior, execution quality, and information asymmetry without the benchmark assumptions required in traditional equity microstructure. I construct seven weekly indices from 233~million resolved Polymarket trades spanning November~2022 to the present, covering probability weighting, market calibration, algorithmic share of volume, longshot-favorite pricing, cross-class execution quality, market efficiency, and wallet-level signatures of private information. The Della Vedova Prediction Market Indices (DV-PMI) are updated weekly, released as CSV panels and per-index JSON snapshots under CC~BY~4.0, and tagged with a semantic version at each refresh. This note documents the construction of each index, reports validation against the underlying working papers, and describes the reproducibility guarantees built into the release process.
\end{abstract}

\section{Introduction}
\label{sec:intro}

Traditional microstructure research relies on fair-value benchmarks that depend on the observed trade sequence, which introduces mechanical correlation between the accuracy and execution components of realized profit \citep{fama1972components,anand2012performance}. Binary prediction markets solve this problem by settling to a known scalar ($\$0$ or $\$1$), removing the need for a modeled benchmark. This paper documents the Della Vedova Prediction Market Indices (DV-PMI), a weekly-updated family of seven behavioral and microstructure indices derived from on-chain Polymarket trade data.

The design follows three conventions. First, each index derives from a working paper whose peer-reviewed analogues predate this release, so the methodological primitives are not novel to DV-PMI. Second, the release cadence is weekly, the schema is pre-registered, and every breaking change is logged in a version-controlled changelog; this matches the pre-registration discipline documented by the Baker-Wurgler sentiment program and the Economic Policy Uncertainty framework. Third, the dataset is open and released under a Creative Commons Attribution 4.0 license for data and an MIT license for code.

The intended uses are three-fold: (i) an external dataset for academic research that requires granular prediction-market state without the engineering cost of independent on-chain extraction; (ii) a citable source for journalists writing about prediction-market activity, regulatory developments, or behavioral bias; and (iii) a methodological reference for expert testimony in regulatory or litigation contexts involving prediction-market trading.

The remainder of this note is organized as follows. Section~\ref{sec:data} describes the on-chain data pipeline and the wallet classification rule. Section~\ref{sec:indices} defines each of the seven indices and its generation procedure. Section~\ref{sec:validation} reports validation against the underlying working papers and out-of-sample checks. Section~\ref{sec:limits} records known limitations. Section~\ref{sec:reproducibility} describes the reproducibility infrastructure. Section~\ref{sec:cite} records the citation format.

\section{Data}
\label{sec:data}

\subsection{Source}
Polymarket operates as a smart-contract binary-outcome exchange on the Polygon network. Every trade is recorded on-chain as an \texttt{OrderFilled} event emitted by one of two exchange contracts (CTF Exchange at \texttt{0x4bFb...82E} and NegRisk Exchange at \texttt{0xC5d5...80a}). The event log contains the maker address, taker address, token identifiers, amount traded, and block number; price is derived from the USDC-denominated numerator and token-denominated denominator of the fill.

Events are collected via \texttt{eth\_getLogs} queries against a Polygon archive node. The collector script retrieves blocks in 2{,}000-block batches, streams events to disk, and applies the \citet{paradigm2023orderfilled} deduplication rule that keeps the maker-side event for each matched trade. Trades with non-positive or above-unity prices are filtered as non-fills. Token identifiers are mapped to market questions and outcome labels via the Polymarket Gamma API.

As of the most recent release, the processed trade panel contains approximately 233 million resolved trades executed by 1.98 million unique wallets across approximately 400{,}000 markets. \TODO{replace placeholders with the refreshed counts as of the as-of date; the current dashboard run reports counts in the latest-snapshot JSON files}.

\subsection{Wallet classification}
Wallets are partitioned into five exclusive classes using the rule from \citet{dellavedova2026profits}:

\begin{itemize}
  \item \textbf{Bot}: more than 50 trades per active day OR more than 1{,}000 total trades.
  \item \textbf{Sophisticated}: lifetime USDC volume above \$10{,}000 AND Herfindahl concentration across markets below 0.5 AND active span above 30 days (and not already classified as Bot).
  \item \textbf{Active Retail}: between 10 and 1{,}000 lifetime trades (and not already classified).
  \item \textbf{Casual}: between 2 and 9 lifetime trades.
  \item \textbf{One-Shot}: exactly one lifetime trade.
\end{itemize}

The classification is applied cumulatively: a wallet's class is determined once per weekly refresh using all trades to date.

\section{Indices}
\label{sec:indices}

All indices are published at weekly resolution, timestamped by ISO week (Monday reference), with a minimum of 13 observations required before reporting a 52-week rolling z-score.

\subsection{Probability Weighting Index (PWI)}
\label{sub:pwi}

Let $\alpha_{c,t}$ denote the one-parameter \citet{prelec1998compound} weighting coefficient fit to the calibration curve of wallet class $c$ in week $t$ by weighted nonlinear least squares, with trade counts as weights and the Prelec form
\[
  w(p;\alpha) = \exp\!\big(-\,(-\log p)^{\alpha}\big).
\]
The PWI is the trade-count-weighted mean of $\alpha_{c,t}$ across all non-bot classes $c$:
\[
  \mathrm{PWI}_t \;=\; \frac{\sum_{c \in \mathcal{C}_{\mathrm{non\text{-}bot}}} n_{c,t}\,\alpha_{c,t}}
                           {\sum_{c \in \mathcal{C}_{\mathrm{non\text{-}bot}}} n_{c,t}},
\]
where $n_{c,t}$ is the trade count of class $c$ in week $t$ and $\mathcal{C}_{\mathrm{non\text{-}bot}} = \{\text{Sophisticated}, \text{Active Retail}, \text{Casual}, \text{One-Shot}\}$. Bots are excluded so the index reflects human probability weighting; bots are best understood as liquidity providers whose implied weighting does not reflect a subjective probability distortion.

Interpretation follows the behavioral literature: $\alpha = 1$ corresponds to the identity ($w(p) = p$, no probability distortion), $\alpha < 1$ to the classical inverse-S pattern documented by \citet{tverskykahneman1992}, and $\alpha > 1$ to the rarer opposite pattern. The pooled non-bot sample matches the \citet{tverskykahneman1992} median estimate of $0.65$ almost exactly, which is reassuring for external validity.

\subsection{Market Calibration Curve}
\label{sub:calibration}

The calibration curve is a cross-sectional snapshot (not a time series). Non-bot trades are pooled across all resolved markets and grouped into 20 price bins of width 5\%. For each bin, the realized win rate is computed along with a 95\% Wilson confidence interval. A one-parameter Prelec function is fit to the 20 bin centers by weighted nonlinear least squares, with trade counts as weights. Reported quantities: $\alpha$, fit $R^2$, per-bin realized frequency, CI, and the fit curve on a 97-point grid.

\subsection{Execution Edge Monitor}
\label{sub:execution}

The Execution Edge is defined as the weekly gap
\[
  \mathrm{Exec}_t \;=\; \alpha_{\mathrm{bot}, t} \;-\; \alpha_{\mathrm{active\text{-}retail}, t}.
\]
A behavioral proxy for the execution-quality divergence documented in \citet{dellavedova2026profits}, where Bots capture $-4.25$~basis-point effective spreads while Active Retail pays $+11.68$~bp. The gap varies over time and tracks both the persistence of that divergence and regime shifts in market composition.

\subsection{Private Information Index (PII)}
\label{sub:pii}

For each wallet $i$, excess accuracy is defined as
\[
  \mathrm{EA}_i \;=\; \mathrm{accuracy}_i \;-\; \overline{\max(p, 1-p)}_i,
\]
where the second term is the accuracy achievable by mechanically following the price signal (picking the side above 50\% in each trade). The null hypothesis is $\mathrm{EA}_i = 0$. A one-sided binomial $z$-test at threshold $z > 2.326$ ($p < 0.01$) with positive $\mathrm{EA}$ flags wallets whose accuracy cannot be explained by price-following alone. Holm-Bonferroni and Benjamini-Hochberg corrections are applied separately. The theta-epsilon orthogonality violation (a positive correlation between the accuracy and execution components at the wallet level, which is absent in the population-level orthogonality established by \citet{dellavedova2026profits}) distinguishes informed traders from lucky ones.

Sub-populations are split by the MNPI taxonomy of \citet{dellavedova2026orthogonality}: \emph{vote} markets (elections, awards), \emph{action} markets (policy decisions, launches), \emph{performance} markets (sports, competitions), and \emph{stochastic} markets (crypto prices, aggregate-force outcomes). The empirical flag-rate ordering matches the a~priori MNPI-risk ordering.

\subsection{Bot Share of Volume}
\label{sub:botshare}

Weekly share of trade count attributable to wallets in the Bot class:
\[
  \mathrm{BotShare}_t \;=\; \frac{n_{\mathrm{bot}, t}}{\sum_c n_{c,t}}.
\]
Headline measure of how algorithmic the market is at any given week.

\subsection{Longshot~/~Favorite Price Gap}
\label{sub:pricegap}

Realized win rate in the longshot price band (prices below 10\%) minus the band midpoint (5\%):
\[
  \mathrm{LongshotGap}_t \;=\; \mathrm{winrate}^{\mathrm{longshot}}_t \;-\; 0.05.
\]
Positive values indicate longshots are under-priced relative to reality; negative values indicate the classical favorite-longshot bias. A favorite-side analog is planned for a future release.

\subsection{Market Efficiency Trend}
\label{sub:efficiency}

Weekly $R^2$ of the one-parameter Prelec fit applied to Active Retail trades:
\[
  \mathrm{Efficiency}_t \;=\; R^2\!\big(\mathrm{Prelec}(\,\cdot\,; \alpha_{\mathrm{active\text{-}retail},t})\big).
\]
Higher values indicate that the week's calibration gap is well-described by a single behavioral parameter; lower values indicate that idiosyncratic shocks dominate the curve. The Active Retail class is chosen because it is the largest human wallet cohort and the most representative of behavioral pricing.

\subsection{Rolling statistics}
Each weekly index is published with 4-, 13-, and 52-week moving averages and a 52-week rolling $z$-score. The $z$-score requires at least 13 non-null observations before reporting and is blank before the 13th week of the series. The landing-page ``notable movements'' block fires when the absolute $z$-score exceeds 1.0.

\section{Validation}
\label{sec:validation}

\subsection{Comparison to working-paper results}

The PWI non-bot pooled $\alpha$ estimate matches the \citet{tverskykahneman1992} experimental median of $0.65$ to within 0.01 under the current sample window, providing face validity for the aggregation. The PII flag count of 6{,}032 wallets reproduces the \citet{dellavedova2026orthogonality} headline count exactly; the Holm-Bonferroni survivor count (746) and BH-FDR survivor count (2{,}337) are likewise reproduced. The MNPI flag-rate ordering (vote~$>$~stochastic~$\approx$~action~$>$~performance) matches the theoretical prediction. The Execution Gap aligns in sign and approximate magnitude with the bps-denominated effective-spread gap documented in \citet{dellavedova2026profits}. \TODO{add a validation table with the precise numbers once the refreshed data lands}.

\subsection{Out-of-sample checks}

\TODO{once the data refresh completes, add a table showing each index's value at the end of the original sample (2026-W05) and at the current as-of date, with a note on stability}.

\subsection{Bot-exclusion diagnostic}

To verify that the PWI's bot exclusion is load-bearing, I compute the bot-inclusive and bot-exclusive variants of the weekly weighted $\alpha$ and report both in the supplementary appendix. The variants differ by roughly 0.4 in weeks with heavy bot volume share and by less than 0.1 in weeks with balanced composition; the exclusion is material.

\section{Known Limitations}
\label{sec:limits}

\textbf{Platform coverage}. DV-PMI v1 covers Polymarket only. The methodology applies to any binary prediction market with public settlement data. Kalshi integration is contingent on wallet-level data access and is planned for a future release.

\textbf{Refresh cadence}. Data refreshes weekly (Sunday 22:00 Pacific). Intraweek changes are not captured. For use cases requiring sub-weekly resolution, users should derive their own aggregations from the raw processed-trade panel.

\textbf{Sample period}. The series begins in November 2022, when Polymarket's on-chain exchange became operational. Pre-2022 prediction-market history (Intrade, PredictIt) is not covered.

\textbf{Wallet identity}. Wallets are cryptographic addresses, not individuals. A single person may operate multiple wallets, and a single wallet may be operated by multiple individuals. The classification rule is a behavioral proxy, not an identity claim.

\textbf{Terminal-payoff assumption}. All indices assume binary settlement at $\{0, 1\}$. Markets with non-binary outcomes, partial resolution, or ambiguous resolution are excluded.

\section{Reproducibility}
\label{sec:reproducibility}

All aggregation code is released under an MIT license in the project GitHub repository (\url{https://github.com/jdellavedova/polymarket-index}). Python dependencies are pinned in \texttt{requirements.txt}. Each weekly refresh is tagged with a semantic version (\texttt{vMAJOR.MINOR.PATCH}) and archived to Zenodo with a permanent DOI. The \texttt{CHANGELOG.md} records every methodology change with date and rationale.

Any change to a published CSV column name, type, or units is a breaking change. It triggers a major version bump with a one-release deprecation window in which the prior column remains alongside its replacement.

The weekly pipeline is:
\begin{enumerate}
  \item \texttt{python polymarket\_blockchain\_collector.py full} extends the raw-events file from the last checkpointed block to the current Polygon tip.
  \item \texttt{python process\_pipeline.py} deduplicates and writes \texttt{processed\_trades.csv}.
  \item A Paper~4-side aggregation step produces the weekly per-type alpha table.
  \item \texttt{python pipeline/run\_all.py} reads the aggregation outputs and writes all DV-PMI CSVs and JSONs.
  \item \texttt{git commit}, \texttt{git tag}, \texttt{git push} deploys via GitHub Actions; Zenodo archives on tag.
\end{enumerate}

\section{Citation}
\label{sec:cite}

Suggested citation:

\begin{quote}
Della Vedova, J. (2026). \emph{Della Vedova Prediction Market Indices}. Version 0.1.0. University of San Diego. \url{https://jdellavedova.com}.
\end{quote}

BibTeX:
\begin{verbatim}
@techreport{dellavedova2026dvpmi,
  title       = {Della Vedova Prediction Market Indices: Construction,
                 Validation, and Use},
  author      = {Della Vedova, Joshua},
  year        = {2026},
  institution = {University of San Diego, Knauss School of Business},
  version     = {0.1.0},
  url         = {https://jdellavedova.com},
  note        = {Weekly-updated index release; data and methodology at
                 https://github.com/jdellavedova/polymarket-index}
}
\end{verbatim}

\bibliographystyle{chicago}
\bibliography{methodology}

\end{document}
